La replicazione dell'infrastruttura di una challenge per consentire la
verifica della patch pone una sfida architetturale fondamentale: i
container clonati devono preservare fedelmente la propria configurazione
di rete (inclusi indirizzi IP statici, interfaccia
\passthrough{\lstinline!eth0!} e default gateway), garantendo
l'esecuzione trasparente degli exploit e dei controlli dell'Oracolo
senza generare conflitti di instradamento o collisioni ARP sull'host con
la rete di produzione già attiva. Al contempo, per supportare scenari
realistici in cui i servizi o gli oracoli necessitano di scaricare
dipendenze o contattare repository esterni, l'ambiente di sandbox deve
poter offrire la connettività WAN verso l'esterno.

Prima di convergere sulla soluzione adottata in BlueAgent --- basata su
una Software-Defined Network orchestrata con Open Virtual Network (OVN)
e un doppio stadio di NAT (illustrata nel Capitolo 4) --- sono state
progettate, implementate e sperimentate quattro architetture di rete
alternative.

In questo capitolo vengono analizzate le caratteristiche tecniche, i
vantaggi e i limiti operativi che hanno portato a scartare ciascuna di
queste soluzioni preliminari.
